\begin{enumerate}
    \item[\textbf{Problem 4.1.}] Show that eq. (4.7) can be obtained by substituting $ix$ for $x$ in eq. (4.6).
    \item[\textbf{Problem 4.2.}] Determine the first four terms of the Taylor expansions of $\tan x$ and $\cot x$ about $x = 0$.
    \item[\textbf{Problem 4.3.}] Determine the first four terms of the Taylor expansions of $\tanh x$ and $\coth x$ about $x = 0$.
    \item[\textbf{Problem 4.4.}] Use dimensional analysis to determine how the catenary parameter, $c$, is related to the constant horizontal component of the cable tension, $T_0$, and its weight per unit length (or unit weight), $\gamma$.
    \item[\textbf{Problem 4.5.}] How much tape sag is permissible to measure a 50 m distance accurately to within 5%? Within 2%?.
    \item[\textbf{Problem 4.6.}] What does a body that weighs 10 N at the earth’s surface weigh at a height of 10 m? At the peak of Mt. Everest? (Hint: You might have to look up some facts about our planet!).
    \item[\textbf{Problem 4.7.}] According to eq. (4.30), at what altitude would the weight of 10 N at the earth’s surface drop to 9 N? To 5 N?.
    \item[\textbf{Problem 4.8.}] Compare the results obtained in Problem 4.7 with more exact results obtained by using eq. (4.29).
    \item[\textbf{Problem 4.9.}] What does a body that weighs 10 N at the earth’s surface weigh on the surface of the moon? On the surface of the planet Pluto? On the surface of the planet Mars? (Hint: You might have to look up some facts about our planet’s environment!).
    \item[\textbf{Problem 4.10.}] If the gravitational potential corresponding to Newton’s law of gravitation (eq. (4.25)) is given by $V_g = -Gm_em / R$, find the exact expression that defines this potential as a function of altitude, $z$, from the earth’s surface.
    \item[\textbf{Problem 4.11.}] Write a binomial expansion of the results of Problem 4.10 to determine the potential energy above the earth’s surface to the first order in $z$.
    \item[\textbf{Problem 4.12.}] Fill in the missing elements of the following table to two-term order: $\sin x$, $\cos x$, $1-\sin x$, $1-\cos x$.
    \item[\textbf{Problem 4.13.}] Fill in the missing elements of the following table to two-term order: $\sinh x$, $\cosh x$, $1-\sinh x$, $1-\cosh x$.
    \item[\textbf{Problem 4.14.}] Develop a volume coefficient of expansion, $\beta$, for a solid of length $L_0$, width $W_0$, and height $H_0$, that parallels the surface coefficient, $\gamma$, of eq. (4.37).
    \item[\textbf{Problem 4.15.}] To what temperature difference would an aluminum solid have to be subjected for the surface coefficient of expansion to produce errors of 1% in the area change compared to the exact area change?.
    \item[\textbf{Problem 4.16.}] To what temperature difference would an aluminum solid have to be subjected for the volume coefficient of expansion to produce errors of 1% in the area change compared to the exact area change?.
    \item[\textbf{Problem 4.17.}] Round off each of the following numbers to two (2) significant figures: (a) 5.237 (b) 0.82549 (c) 81.356 (d) $\pi$ (e) 6.2305 (f) 0.0428 (g) 10.45 (h) 4.035.
    \item[\textbf{Problem 4.18.}] Round off each of the following numbers to three (3) significant figures: (a) 5.237 (b) 0.82549 (c) 81.356 (d) $\pi$ (e) 6.2305 (f) 0.0428 (g) 10.45 (h) 4.035.
    \item[\textbf{Problem 4.19.}] Complete the following multiplications and express the results to the correct number of significant figures: (a) $(6.28 \times 10^3) \times 2.712$ (b) $43.32 \times 0.3$ (c) $928 \times 4.23$.
    \item[\textbf{Problem 4.20.}] Do 99.9 and 100.1 have the same number of significant figures? Explain your answer.
    \item[\textbf{Problem 4.21.}] Estimate the ranges within which each of the following numbers lie: (a) 7.7 (b) 7.70 (c) 1200 (d) $1.200 \times 10^{-3}$.
    \item[\textbf{Problem 4.22.}] By what percentage would the period of a pendulum change if its length was halved? If it was reduced by one-third? If the length was reduced to one-third of its original length?.
    \item[\textbf{Problem 4.23.}] Explain why the pendulum period increases by 145% on the moon.
    \item[\textbf{Problem 4.24.}] How would the period of a pendulum change, compared to its value on earth, if the pendulum was on Mars? On Pluto?.
    \item[\textbf{Problem 4.25.}] How would the period of a pendulum change as a function of its height, $h$, above the surface of the earth? (Hint: The variation of the gravitational acceleration $g$ can be represented as a function of $h$ from Newton’s law of gravitational attraction).
    \item[\textbf{Problem 4.26.}] Draw two circular archery targets and use them to depict the hit patterns of (a) an archer who is accurate, but not precise; and (b) an archer who is precise, but not accurate.
    \item[\textbf{Problem 4.27.}] Verify the final forms of eqs. (4.51) and (4.52).
    \item[\textbf{Problem 4.28.}] Verify the equations for $m$ and $b$ given in eqs. (4.53) and (4.54).
    \item[\textbf{Problem 4.29.}] Discuss and explain the dimensional differences between eqs. (4.53) and (4.54).
    \item[\textbf{Problem 4.30.}] Verify the terms in the third and fourth columns of Table 4.3, as well as the sums of all four columns.
    \item[\textbf{Problem 4.31.}] Verify the calculations of $m$ and $b$ found from the results in Table 4.3.
    \item[\textbf{Problem 4.32.}] Determine the standard deviation for the data presented in Table 4.4.
    \item[\textbf{Problem 4.33.}] Draw a histogram for the data in Table 4.4 with 10 intervals of 4 dB width.
    \item[\textbf{Problem 4.34.}] Show that the square of the sample variance of eq. (4.57) can be cast in the alternative form $s^2 = \frac{1}{(n-1)} \left[ \sum_{i=1}^n x_i^2 - \frac{1}{n} \left(\sum_{i=1}^n x_i\right)^2 \right]$.
    \item[\textbf{Problem 4.35.}] Estimate the error made in approximating $y(x) = \sin x$ with a Taylor’s formula to $n = 4$ by evaluating the remainder $R_5$.
    \item[\textbf{Problem 4.36.}] Do the statements that $\sin x \simeq x$ and $\tan x \simeq x$ produce similar approximations? Confirm and explain your answer.
    \item[\textbf{Problem 4.37.}] The readings of an old-fashioned analog voltmeter are subject to some systematic error where all of its readings are too large. The magnitude of the error has been found to vary linearly from 1 V at a dial reading of 5 V to 4 V at a dial reading of 80 V. (a) What are the correct voltages for dial readings of 80, 100, 50, 1, 35, and 10 V? (b) What is the percentage error for each of the six readings in part (a)?.
    \item[\textbf{Problem 4.38.}] (a) Is it possible to have a set of measurements that are precise but not accurate? Explain. (b) Is it possible to have a set of measurements that are accurate but not precise? Explain.
    \item[\textbf{Problem 4.39.}] (a) Write the Taylor series expansion for $e^x$ about $x = 0$. (b) Calculate $e^{0.5}$ to five significant figures using the first four terms of the series found in part (a).
    \item[\textbf{Problem 4.40.}] (a) What percentage error was incurred in the calculation of part (b) of Problem 4.39 if the true value of $e^{0.5}$ is 1.6487? (b) Use the Taylor remainder (eq. (4.5)) to calculate the error in $e^{0.5}$ after only four terms. Is the error calculated in part (a) acceptable? Explain.
    \item[\textbf{Problem 4.41.}] Evaluate the following function by hand for $x = 4$: $(1 + 2/x)^{1/4}$.
    \item[\textbf{Problem 4.42.}] How does an observer know when enough measurements have been taken?.
    \item[\textbf{Problem 4.43.}] Make a list of five new examples of systematic errors.
    \item[\textbf{Problem 4.44.}] Make a list of five new examples of random errors.
    \item[\textbf{Problem 4.45.}] The resistance of a resistor, $R$, is estimated by passing several currents, $I$, through it and measuring corresponding voltage drops, $V$, and currents with imprecise analog meters. Data: $x_i = V$ (V): 10, 20, 30, 40, 50, 60, 70, 80; $y_i = I$ (A): 0.8, 1.1, 2.5, 4.2, 4.3, 4.7, 5.8, 6.4. (a) What kinds of error will be found? (b) Assuming $V = IR$, plot and eyeball the best-fit line.
    \item[\textbf{Problem 4.46.}] Use the method of least squares to plot a $V$ versus $I$ curve for the data of Problem 4.45. How does it compare with the eyeball result?.
    \item[\textbf{Problem 4.47.}] The data comprise 100 readings of noise levels taken 6 mi away from an airport, at 15 s intervals. Find mean, median, and standard deviation.
    \item[\textbf{Problem 4.48.}] The starred numbers in Problem 4.47 were taken while an aircraft flew overhead. If those are deleted, what are the mean, median, and standard deviation of the remaining 88 points?.
    \item[\textbf{Problem 4.49.}] Draw (a) a histogram of all data of Problem 4.47 and (b) a continuous curve of number of readings versus measured noise level.
    \item[\textbf{Problem 4.50.}] Determine a far-field approximation of $f(r)=\sqrt{a^2+r^2}$ as a binomial expansion for $r \gg a$.
    \item[\textbf{Problem 4.51.}] The electric potential at distance $r$ along the axis of a disk of radius $a$ is $V_e = \frac{q}{2\pi a^2 \varepsilon_0}(\sqrt{a^2 + r^2} - r)$. Using Problem 4.50, find a far-field approximation for $r \gg a$.
    \item[\textbf{Problem 4.52.}] Compare the minimum number of terms kept in the binomial expansions of Problems 4.50 and 4.51. Are they the same? Why or why not?.
    \item[\textbf{Problem 4.53.}] Suppose we need radial extension or deflection $w$ of a very thin spherical balloon, with radius extending from $R$ to $R+w$ under pressure. Formula proposed: $w/R = pR/Eh$, where $h$ is wall thickness and $E$ is modulus. Is this dimensionally consistent?.
    \item[\textbf{Problem 4.54.}] Analyze limit behavior of Problem 4.53 as pressure, modulus, radius, and thickness go to zero and to infinity. Does this behavior match intuition?.
    \item[\textbf{Problem 4.55.}] Use the equation in Problem 4.53 to estimate pressure magnitude $p$ as a fraction of modulus $E$. Estimate this fraction for a thin-walled sphere with $h/R \ll 1$.
\end{enumerate}
